\subsection*{Abstract}

Phylodynamic models use phylogenetic branching patterns to infer population and diversification processes, including speciation, extinction, and epidemiological transmission.
Extensions of these models enable the incorporation of external covariates as predictors of phylodynamic parameters, for instance allowing us to link traits or environmental variables with changes in speciation, extinction, or transmission rates. 
However, the dependencies between predictors and phylodynamic parameters are typically restricted to linear, additive effects and may thus fail to capture complex, non-linear relationships underlying evolutionary dynamics.
To address this limitation, we propose a new framework that leverages unsupervised Bayesian neural networks to flexibly model functional relationships between phylodynamic parameters and a broad set of predictors, including categorical traits, quantitative variables, and time series data.
Based on these covariates, the network weights are estimated jointly with the phylogenetic tree, obtained directly from sequence alignment data. Using extensive simulations, we demonstrate that this approach accurately recovers predictor-parameter relationships, mitigates overfitting, and remains robust across both macroevolutionary and epidemiological contexts.
By incorporating tools from explainable artificial intelligence, we further show that our framework reliably identifies the most influential predictors and yields interpretable descriptions of their impact on phylodynamic rates.
Finally, we apply our method to two empirical analyses: linking SARS-CoV-2 migration dynamics with travel data during its early spread in Europe, and inferring trait and time-dependent speciation and extinction rates in the diversification of platyrrhines. Our framework substantially expands the capabilities of phylodynamic inference providing a powerful and flexible approach to model complex macroevolutionary and epidemiological processes.
